% ============================================================
% Section 56: 格林艾森常量与晶格热膨胀
% ============================================================
\newpage
\section{固体理论：格林艾森常量与晶格热膨胀}
\label{sec:gruneisen-thermal-expansion}

\begin{modelbox}[题目：简谐晶体的自由能与热膨胀]
考虑一个内部原子振动的固体。
\begin{enumerate}
    \item 计算温度为 $T$ 时频率为 $\omega$ 的声子模的自由能 $F$；
    \item 假设此固体是简谐的，体积弹性模量为 $K$，相对体积变化为 $\Delta=(V-V_0)/V_0$，忽略声子模的色散关系，即认为 $\omega_q=\omega$，写出晶体的自由能；
    \item 如果 $\omega$ 与体积关系为 $\displaystyle\frac{\delta\omega}{\omega}=-\gamma\frac{\delta V}{V_0}$（$\delta V$ 为体积变化量），其中 $\gamma$ 称为格林艾森常量，那么温度为 $T$ 时相对体积变化 $\Delta$ 为多少？
    \item 讨论格林艾森常量的物理意义。
\end{enumerate}
\end{modelbox}

\begin{tipbox}[考点快速分析]
\begin{itemize}
    \item \textbf{单模自由能}：由配分函数 $Z$ 出发，$F=-k_BT\ln Z$
    \item \textbf{晶体自由能}：静态晶格能 $U(V)$ + 所有声子模自由能之和
    \item \textbf{状态方程}：$P=-(\partial F/\partial V)_T$，热平衡时 $P=0$
    \item \textbf{格林艾森关系}：$\partial\ln\omega/\partial\ln V=-\gamma$，联系振动频率与体积变化
    \item \textbf{热膨胀}：$\Delta=\gamma E/(KV_0)$，非简谐效应的直接体现
\end{itemize}
\end{tipbox}

\begin{derivationbox}[标准解题过程]
\textbf{Step 1: 单个声子模的自由能}

频率为 $\omega$ 的量子谐振子，能级为 $E_n=(n+\tfrac{1}{2})\hbar\omega$（$n=0,1,2,\ldots$）。配分函数：
\begin{equation}
Z=\sum_{n=0}^{\infty}e^{-(n+1/2)\hbar\omega/k_BT}
=e^{-\hbar\omega/2k_BT}\sum_{n=0}^{\infty}\bigl(e^{-\hbar\omega/k_BT}\bigr)^n
=\frac{e^{-\hbar\omega/2k_BT}}{1-e^{-\hbar\omega/k_BT}}.
\end{equation}

自由能 $F=-k_BT\ln Z$：
\begin{align}
F_1&=-k_BT\Bigl[-\frac{\hbar\omega}{2k_BT}-\ln\bigl(1-e^{-\hbar\omega/k_BT}\bigr)\Bigr]\\
&=\frac{1}{2}\hbar\omega+k_BT\ln\bigl(1-e^{-\hbar\omega/k_BT}\bigr).
\end{align}

其中 $\tfrac{1}{2}\hbar\omega$ 为零点能，$k_BT\ln(1-e^{-\hbar\omega/k_BT})$ 为热激发贡献。

\textbf{Step 2: 简谐晶体的自由能}

假设所有振动模式具有相同频率 $\omega$（爱因斯坦近似），总自由度为 $N$（三维单原子晶体 $N=3N_{\text{atom}}$）。晶体自由能为静态晶格能 $U(V)$ 与所有声子模自由能之和：
\begin{equation}
F(V,T)=U(V)+N\Bigl[\frac{1}{2}\hbar\omega+k_BT\ln\bigl(1-e^{-\hbar\omega/k_BT}\bigr)\Bigr].
\label{eq:crystal-free-energy}
\end{equation}

\textit{自由能的物理意义}：$F(V,T)$ 是描述晶体热力学平衡的势函数。给定温度 $T$，晶体自发调整体积 $V$ 使 $F$ 取极小值。它包含两部分：
\begin{itemize}
    \item $U(V)$：\textbf{静态晶格能}（$T=0$ 时的原子相互作用势）。即使温度为零，原子间也存在相互作用（吸引+排斥），这部分能量与热振动无关，只由原子间距（即体积）决定。它决定了晶体的``骨架''——平衡体积 $V_0$ 和结合强度。
    \item 声子项：\textbf{热振动的量子化能量}。声子不是势能本身，而是原子在平衡位置附近振动的\textbf{激发量子}。把原子从静止的平衡位置``推''一下，它以频率 $\omega$ 振动，这就是声子。声子的能量（零点能+热激发能）叠加在静态晶格能之上，构成总自由能。
\end{itemize}

\textit{类比}：把晶体想象成一张拉紧的鼓面。$U(V)$ 是鼓面被拉紧时储存的弹性势能（与绷紧程度有关）；声子项是敲击鼓面后产生的声波（振动模式）的能量。鼓面越紧（$U$ 越大），声波传播越快（$\omega$ 越大），二者通过 $U(V)$ 间接耦合——这就是热膨胀的微观根源。

\textbf{Step 3: 热膨胀（或收缩）的计算}

\textit{状态方程}：压强由自由能对体积的偏导给出：
\begin{equation}
P=-\Bigl(\frac{\partial F}{\partial V}\Bigr)_T
=-\frac{\partial U}{\partial V}-N\Bigl(\frac{1}{2}\hbar+\frac{\hbar}{e^{\hbar\omega/k_BT}-1}\Bigr)\frac{\partial\omega}{\partial V}.
\end{equation}

括号内的项正是单个模式的平均振动能：
\begin{equation}
\bar{E}_1=\frac{1}{2}\hbar\omega+\frac{\hbar\omega}{e^{\hbar\omega/k_BT}-1}.
\end{equation}

总振动能 $E=N\bar{E}_1$，状态方程可简洁写为
\begin{equation}
P=-\frac{\partial U}{\partial V}-\frac{E}{\omega}\frac{\partial\omega}{\partial V}.
\label{eq:state-equation}
\end{equation}

\textit{格林艾森关系}：格林艾森常量的严格定义是
\begin{equation}
\gamma=-\frac{d\ln\omega}{d\ln V}=-\frac{V}{\omega}\frac{d\omega}{dV}.
\label{eq:gruneisen-relation}
\end{equation}

它描述体积的无穷小相对变化 $dV/V$ 引起频率的相对变化 $d\omega/\omega$。对有限小变形，一阶泰勒展开给出
\begin{equation}
\frac{\delta\omega}{\omega}\approx-\gamma\frac{\delta V}{V_0}=-\gamma\Delta,
\qquad\text{即}\qquad
\frac{\partial\omega}{\partial V}\approx-\gamma\frac{\omega}{V_0}.
\end{equation}

\textit{符号约定}：$\Delta=(V-V_0)/V_0$ 为相对体积变化；$\delta V=V-V_0$ 为体积变化量。$\gamma>0$ 时，体积增大（$\Delta>0$）导致频率降低（$\delta\omega<0$）——原子间距变大，弹簧变``松''，振动变慢。

代入 \eqref{eq:state-equation}：
\begin{equation}
P=-\frac{\partial U}{\partial V}+\gamma\frac{E}{V}.
\end{equation}

\textit{热平衡条件}：无外压时 $P=0$，晶体体积由热振动自发调整：
\begin{equation}
\frac{\partial U}{\partial V}=\gamma\frac{E}{V}.
\end{equation}

将 $\partial U/\partial V$ 在静态体积 $V_0$ 附近泰勒展开至一阶。令 $\delta V=V-V_0$，则
\begin{equation}
\frac{\partial U}{\partial V}=\Bigl(\frac{\partial U}{\partial V}\Bigr)_{V_0}+\Bigl(\frac{\partial^2U}{\partial V^2}\Bigr)_{V_0}\delta V.
\end{equation}

静态平衡时 $(\partial U/\partial V)_{V_0}=0$。体弹性模量定义为 $K=V_0(\partial^2U/\partial V^2)_{V_0}$，故
\begin{equation}
\frac{\partial U}{\partial V}\approx\frac{K}{V_0}\delta V=K\Delta.
\end{equation}

代入热平衡条件（小变形下 $V\approx V_0$）：
\begin{equation}
K\Delta=\gamma\frac{E}{V_0}
\quad\Longrightarrow\quad
\boxed{\Delta=\frac{\gamma E}{KV_0}}.
\end{equation}

由于 $E>0$，若 $\gamma>0$ 则 $\Delta>0$，表现为\textbf{热膨胀}（$V>V_0$）；若 $E$ 小于零温参照值（如负温度情形），则相应为\textbf{收缩}（$V<V_0$）。通常情况下热振动能随温度升高而增大，故 $\Delta>0$。

\textbf{Step 4: 格林艾森常量的物理意义}

\textit{定义}：
\begin{equation}
\gamma=-\frac{d\ln\omega}{d\ln V}
=-\frac{V}{\omega}\frac{d\omega}{dV}.
\end{equation}

用 $d$ 而非 $\partial$ 是因为在爱因斯坦近似下，$\omega$ 仅是体积 $V$ 的单变量函数，与温度 $T$ 无关（温度只改变声子占据数，不改变频率本身）。

\textit{与热膨胀的关系}：由 $\Delta=\gamma E/(KV_0)$ 对 $T$ 求导，得热膨胀系数
\begin{equation}
\alpha=\frac{1}{V_0}\frac{\partial(\delta V)}{\partial T}=\frac{\partial\Delta}{\partial T}
=\frac{\gamma}{KV_0}\frac{\partial E}{\partial T}
=\frac{\gamma C_V}{KV_0}.
\end{equation}

\textit{物理意义总结}：
\begin{itemize}
    \item \textbf{非简谐效应的度量}：严格简谐晶体中振动频率与体积无关，$\gamma=0$，则热膨胀系数 $\alpha=0$。$\gamma\neq0$ 是晶格非简谐性的直接体现。
    \item \textbf{``简谐''与``非简谐''的辩证}：题目说``固体是简谐的''，是指原子在平衡位置附近的振动可用简谐近似处理（势能展开到二阶）；但同时又允许 $\omega$ 依赖于 $V$（通过格林艾森关系），这种依赖实际上是势函数三阶及以上项（非简谐项）的结果。这种处理方式称为\textbf{准简谐近似}（quasi-harmonic approximation）——振动本身是简谐的，但振动参数（频率）由非简谐效应决定。它是计算有限温度晶格性质的标准方法。
    \item \textbf{热学与力学的桥梁}：格林艾森常量将固体的热学性质（热容、热膨胀）与力学性质（弹性模量、状态方程）定量联系起来，是 Mie-Gr\"uneisen 物态方程的核心参数。
    \item \textbf{微观图像}：$\gamma$ 的大小反映原子间相互作用势的非对称性——排斥势随距离变化陡峭，吸引势相对平缓。典型固体的 $\gamma$ 值在 $1\sim3$ 之间。
\end{itemize}
\end{derivationbox}

\begin{formulabox}[答案汇总]
\begin{center}
\begin{tabular}{cl}
\toprule
问题 & 结果 \\
\midrule
(1) 单模自由能 & $F_1=\dfrac{1}{2}\hbar\omega+k_BT\ln\bigl(1-e^{-\hbar\omega/k_BT}\bigr)$ \\[10pt]
(2) 晶体自由能 & $F=U(V)+N\bigl[\tfrac{1}{2}\hbar\omega+k_BT\ln(1-e^{-\hbar\omega/k_BT})\bigr]$ \\[10pt]
(3) 热膨胀 & $\Delta=\dfrac{\gamma E}{KV_0}$ \\[10pt]
(4) $\gamma$ 的物理意义 & 非简谐效应度量；$\gamma=0$ 则无热膨胀 \\
\bottomrule
\end{tabular}
\end{center}
\end{formulabox}

\begin{insightbox}[物理图像与概念深化]
\textbf{1. 为什么简谐晶体没有热膨胀？}

简谐近似下，原子间相互作用势为抛物线 $U(r)\approx U_0+\tfrac{1}{2}\beta(r-r_0)^2$。势能关于平衡位置对称：原子在平衡位置两侧的概率分布相同，平均位移为零。无论温度多高，原子振动始终对称地分布在 $r_0$ 两侧，不会导致平均间距变化。

热膨胀来源于势能曲线的\textbf{非对称性}（非简谐效应）：
\begin{itemize}
    \item 排斥支陡峭（短程强斥力）
    \item 吸引支平缓（长程弱引力）
    \item 温度升高时，原子更倾向于向排斥侧偏移（因为吸引侧``更深更宽''），平均间距增大
\end{itemize}

格林艾森常量 $\gamma$ 正是对这种非对称性的定量刻画。

\textbf{2. 自由能的两部分贡献}

\eqref{eq:crystal-free-energy} 中：
\begin{itemize}
    \item $U(V)$：静态晶格能（$T=0$ 时的原子相互作用势），决定晶体的平衡结构和体积
    \item 声子项：热振动的贡献，包含零点能（量子效应）和热激发能（温度依赖）
\end{itemize}

在热力学中，$F(V,T)$ 是描述晶体平衡性质的势函数：给定 $T$，晶体自发调整 $V$ 使 $F$ 最小。这正是热膨胀的微观机制。

\textbf{3. 格林艾森常量的典型数值}

\begin{center}
\begin{tabular}{cccc}
\toprule
材料 & $\gamma$ & 材料 & $\gamma$ \\
\midrule
Cu & $\sim1.9$ & NaCl & $\sim1.5$ \\
Al & $\sim2.2$ & KCl & $\sim1.4$ \\
Fe & $\sim1.6$ & 金刚石 & $\sim1.1$ \\
\bottomrule
\end{tabular}
\end{center}

金刚石 $\gamma$ 较小（$\sim0$），说明其共价键势阱接近对称，但仍存在微弱的非简谐性。金刚石的低热膨胀系数（室温下 $\sim10^{-6}\,\mathrm{K^{-1}}$）与其高硬度、高德拜温度一致。

\textbf{4. 与德拜模型的关联}

\ref{sec:debye-phonon-number}（德拜模型）中假设频率与体积无关，即 $\partial\omega/\partial V=0$，$\gamma=0$。这是\textbf{简谐近似}的结果，因此德拜模型本身不能解释热膨胀。要在德拜框架下讨论热膨胀，必须引入格林艾森修正：假设所有模式的 $\omega$ 都以相同的 $\gamma$ 随体积变化，这就是\textbf{德拜-格林艾森模型}。

\textbf{5. 负热膨胀（NTE）材料}

某些材料（如 $\mathrm{ZrW_2O_8}$、部分框架结构材料）在特定温区表现出\textbf{负热膨胀}（温度升高，体积收缩）。其微观机制通常与：
\begin{itemize}
    \item 特殊的声子模式（如转动机理，rattling phonons）
    \item 强的非简谐耦合导致某些模式 $\gamma<0$
    \item 电子贡献（如巡游电子体系的磁性相变）
\end{itemize}

负热膨胀材料在精密仪器（如光学平台、航天器结构）中有重要应用价值。
\end{insightbox}

\begin{thinkbox}[考场速记：格林艾森问题答题框架]
遇到``热膨胀与格林艾森常量''类问题：
\begin{enumerate}
    \item \textbf{写单模自由能}：$F_1=\tfrac{1}{2}\hbar\omega+k_BT\ln(1-e^{-\hbar\omega/k_BT})$
    \item \textbf{写晶体自由能}：$F=U(V)+\sum F_1$
    \item \textbf{状态方程}：$P=-(\partial F/\partial V)_T=0$，引入格林艾森关系 $\partial\omega/\partial V=-\gamma\omega/V$
    \item \textbf{泰勒展开}：$\partial U/\partial V\approx K\Delta$，解得 $\Delta=\gamma E/(KV_0)$
    \item \textbf{物理意义}：$\gamma$ 衡量非简谐性，$\gamma=0$ 无热膨胀
\end{enumerate}
\textit{核心逻辑链}：温度 $\to$ 热振动能 $E$ $\to$ 通过 $\gamma$ 耦合到体积 $V$ $\to$ 弹性恢复力（$K$）抵抗膨胀 $\to$ 平衡时 $\Delta=\gamma E/(KV_0)$。
\end{thinkbox}
